\chapter{الأقفال والتزامن}
\label{chap:lock}

تتطلب معالجة تعدد الأنوية وتعدد العمليات استخدام آليات محكمة للتحكم في الوصول المباشر لبيانات النواة المشاركة؛
فإذا حاولت نواتان تعديل قائمة الذاكرة الحرة في نفس الوقت دون تزامن، فقد تخرب بنية البيانات وتتداخل المؤشرات.

تُسمى الآلية الأساسية للتزامن بـ \indextext{القفل (Lock)}، حيث يضمن القفل غياب التداخل والوصول التبادلي المتبادل (Mutual Exclusion) لبُنى البيانات الحساسة.

تستخدم \texttt{xv6} نوعين أساسيين من الأقفال:
\begin{enumerate}
\item \textbf{أقفال الدوران (Spinlocks)}: ينتظر المعالج المداعي بداخل حلقة تكرارية مستمرة طالما كان القفل مشغولاً. وتناسب أقفال الدوران العمليات السريعة جداً بداخل النواة مع تعطيل المقاطعات على النواة الحالية لمنع التكالب والانسداد المتبادل (Deadlock).
\item \textbf{أقفال النوم (Sleeplocks)}: تتيح للعملية التنازل عن المعالج والنوم إذا كان القفل مشغولاً، مما يسمح لمعالجات أخرى بالعمل وتجنب ضياع دورات المعالج بداخل حلقات الانتظار.
\end{enumerate}

يتطلب التعامل مع الأقفال حذراً شديداً لتجنب \indextext{الانسداد المتبادل (Deadlock)}،
والذي يحدث عندما تنتظر العملية الأولى القفل \texttt{A} وهي تحوز القفل \texttt{B}، بينما تنتظر العملية الثانية القفل \texttt{B} وهي تحوز القفل \texttt{A}.
تفرض النواة نظاماً صارماً يحدد الترتيب العالمي لاستحواذ الأقفال لمنع الانسداد.
